% ===== tr-01 块2：第 2 章 pp.6--10（原文页 6--10）=====
% [原文 p.6]
另一组对 SYZ 镜像对称与特殊拉格朗日几何的发展具有历史性影响的非紧例子由 Joyce [103] 发现。人们曾相信特殊拉格朗日纤维化总是光滑的，因而其判别轨迹具有余维二。但 Joyce 的例子表明情况很可能并非如此。他构造的是显式的 $S^1$-不变特殊拉格朗日纤维化，它们只是分段光滑的，且判别轨迹具有实余维一。这种纤维化的奇点集是一个 Riemann 曲面，其 amoeba 形状的像给出余维一的判别轨迹。Joyce 论证说，他的例子展示了特殊拉格朗日纤维化判别轨迹的典型行为。

Joyce 的这项重要工作不仅加深了我们对特殊拉格朗日纤维化可能奇点的理解，还迫使人们重新思考 SYZ 猜想。最初，人们期望一对镜像 Calabi--Yau 流形 $X$ 与 $\check X$ 具有到同一个底 $B$、并共享同一判别轨迹的特殊拉格朗日环面纤维化。现在 Joyce 的例子表明：一侧的判别轨迹可以具有余维二，而镜像一侧的判别轨迹却完全可能具有余维一。人们所能期望的最好情形是：当逼近大复结构极限时，两侧的判别轨迹收敛到 $B$ 中同一个余维二子集。

更确切地说，设 $\mathcal{X}\to D$ 与 $\check{\mathcal{X}}\to D$ 是互为镜像的 Calabi--Yau 流形的极大幂零退化，其中 $D$ 是单位圆盘，$0\in D$ 是两侧的大复结构极限。取一列 $\{t_i\}\subset D$ 收敛到 $0$，并设 $g_i,\check g_i$ 分别是 $\mathcal{X}_{t_i},\check{\mathcal{X}}_{t_i}$ 上的 Ricci 平坦度量，将它们归一化使直径固定为 $C$。那么 SYZ 猜想的一个极限版本可以表述为：
\begin{enumerate}
\item[(i)] 存在 $(\mathcal{X}_{t_i},g_i)$ 与 $(\check{\mathcal{X}}_{t_i},\check g_i)$ 的收敛子列，分别（在 Gromov--Hausdorff 意义下）收敛到度量空间 $(B_\infty,d_\infty)$ 与 $(\check B_\infty,\check d_\infty)$；
\item[(ii)] 空间 $B_\infty$ 与 $\check B_\infty$ 是带奇点的仿射流形，且二者都同胚于 $S^n$；
\item[(iii)] 在实余维二轨迹 $\Gamma\subset B_\infty$（相应地 $\check\Gamma\subset\check B_\infty$）之外，$d_\infty$（相应地 $\check d_\infty$）由 Monge--Ampère 度量诱导；且
\item[(iv)] Monge--Ampère 流形 $B_\infty\setminus\Gamma$ 与 $\check B_\infty\setminus\check\Gamma$ 互为 Legendre 对偶。
\end{enumerate}
% [Tutor Note: 原文 p.6 从"More precisely"起用花体 $\mathcal{X},\check{\mathcal{X}}$ 记退化的 Calabi--Yau 族（全族），而 $\mathcal{X}_{t_i}$ 是 $t_i$ 处的纤维；不要把 $\mathcal{X}$ 抄成普通斜体 $X$，二者在原文中是不同的符号。]
% [Tutor Note: 上述 (i)--(iv) 是 SYZ 猜想的"极限版本"（由 Gross--Wilson 与 Kontsevich--Soibelman 独立提出）；$d_\infty$ 是 Gromov--Hausdorff 极限度量，不是原来的 Ricci 平坦度量。]
% [Tutor Note: 原文 p.6 末 $\check X_0:=T\check B_0/\Lambda$ 中的 $T\check B_0$ 表示 $\check B_0$ 的切丛（与 p.4 的 $TB/\Lambda$ 同一记法），并非"乘以 $T$"。]

这独立地由 Gross--Wilson [92] 与 Kontsevich--Soibelman [114] 提出。事实上，理解 Ricci 平坦度量极限行为的一般问题由 Yau 在其著名的公开问题清单 [179, 180] 中提出。受镜像对称的 SYZ 图景启发，这个问题在过去 15 年里得到了广泛研究，Gross--Wilson [92]、Tosatti [166, 167]、Ruan--Zhang [154]、Zhang [181]、Rong--Zhang [149, 150]，以及更近的 Gross--Tosatti--Zhang [91, 90] 取得了实质性进展。

度量空间 $B_\infty$ 与 $\check B_\infty$ 应当被看作 Calabi--Yau 族上 SYZ 纤维化之底的极限。应用上述图景，人们可以尝试如下构造 Calabi--Yau 流形极大幂零退化 $\mathcal{X}\to D$ 的镜像：先确定 Gromov--Hausdorff 极限 $B_\infty$，取 $B_\infty\setminus\Gamma$ 的 Legendre 对偶 $\check B_0$，然后紧化商空间 $\check X_0:=T\check B_0/\Lambda$。

不幸的是，这个途径永远行不通，因为 $\check X_0$ 上的半平坦复结构并非整体定义，其原因是 $\check B_\infty$ 中奇点附近的仿射结构具有非平凡单值。要得到正确的镜像，我们需要按照 SYZ 的建议，用来自全纯圆盘的量子修正来形变 $\check X_0$ 的复结构。

% [原文 p.7]
\section{紧 Calabi--Yau 流形的 SYZ}
人们相信，用半平坦 SYZ 镜像对称得到的辛结构可以自然地紧化，从而在镜像上给出整体的辛结构。确实，Castaño-Bernard 与 Matessi [15] 证明了 Gross 在 [72] 中构造的拓扑紧化可以被赋予辛结构，由此得到成对的紧辛 6-流形，它们同胚于已知的镜像 Calabi--Yau 三维流形对（包括五次三维流形及其镜像），并配有拉格朗日环面纤维化，其底是互为 Legendre 对偶的带奇点整仿射流形。
% [Tutor Note: "compact symplectic 6-folds"（紧辛 6-流形）对应实六维、即复三维的 Calabi--Yau 流形；原文在紧化后保留了辛结构但不再是复流形，故称 6-folds 而非 3-folds。]

另一方面，正如上面提到的，原始的 SYZ 提议 [163] 指出：镜像上的 Ricci 平坦度量应当与半平坦 Calabi--Yau 度量 [68] 相差由全纯圆盘带来的瞬子修正，这些圆盘的边界非平凡地缠绕在 SYZ 纤维化的纤维上。由于镜像上的度量由其辛结构与复结构唯一决定，自然可以期望瞬子修正可用来扰动镜像的复结构。这就是 SYZ 猜想背后的关键思想。由于全纯圆盘可以粘接成全纯曲线，这也解释了为什么镜像对称可以用来解决枚举问题。

给定带奇点的仿射流形 $B$，记 $\Gamma\subset B$ 与 $B_0=B\setminus\Gamma$ 分别为其奇异轨迹与光滑轨迹。我们希望构造一个复流形 $X$，作为 $X_0:=TB_0/\Lambda$ 的一个小形变的部分紧化。这在镜像对称中被称为重建问题。
% [Tutor Note: $B_0=B\setminus\Gamma$ 是 $B$ 的光滑轨迹（去掉奇异轨迹 $\Gamma$ 后）；$TB_0$ 是其切丛，故 $X_0=TB_0/\Lambda$ 是半平坦模型。]

这个问题最早由 Fukaya 在 [53] 中研究，他试图通过直接求解控制 $X_0$ 上复结构形变的 Maurer--Cartan 方程来找到合适的扰动。在二维情形，他的启发式论证表明，所需的扰动应当来自 $B$ 中从 $\Gamma$ 发出的梯度流树。后者应当来自当逼近大复结构极限时，界定 SYZ 纤维化的拉格朗日环面纤维的全纯圆盘的极限。Fukaya 提出了一系列引人入胜的猜想，解释量子修正是如何改变镜像复结构的。然而，要把他那直观上清晰的图景变得严格所需的分析似乎遥不可及。\footnote{在最近的一篇工作 [25] 中，人们利用 [31] 所发展的 Witten--Morse 理论与 de Rham 理论之间的关系实现了 Fukaya 的纲领。特别地，这给出散射图的一种新的几何解释。}

Kontsevich--Soibelman [115] 通过在刚性解析空间上工作，绕过了 Fukaya 论证中的分析困难。他们从 $S^2$ 上一个带 24 个奇点的整仿射结构出发（使绕每个奇点的仿射结构单值为 $\begin{pmatrix}1&1\\0&1\end{pmatrix}$），并成功重建了一个非阿基米德解析 K3 曲面。基本想法是：在 Fukaya 的构造中给每条梯度流线附加一个自同构，并用这些自同构修改镜像上坐标卡之间的粘接，从而抵消判别轨迹附近仿射结构的非平凡单值所带来的影响。
% [Tutor Note: 单值矩阵 $\begin{pmatrix}1&1\\0&1\end{pmatrix}$ 是整仿射变换（对角元为 1、右上元为整数 1），它正是半平坦复结构无法整体定义的根源：绕奇点一圈后坐标卡相差一个非平凡整仿射变换。]

% [原文 p.8]
他们论证中的一个关键步骤是一条关键引理：当两条梯度流线相交时，总可以加入带有自同构的新流线，使得绕每个交点的复合为恒等。这被称为散射图；它保证沿一条路径所穿过的各条流线所附的自同构之复合与所选择的路径无关，从而使修改后的粘接是相容的。

大约在同一时期，Gross 与 Siebert 启动了他们的宏大纲领 [84, 85, 86, 73, 87]，旨在为 SYZ 猜想建立代数几何的途径。受上一节讨论的 SYZ 猜想极限版本以及 Kontsevich 的观察（Gromov--Hausdorff 极限大致就是退化的对偶相交复形）的启发，他们提出了一个代数几何的 SYZ 程序来构造镜像。

从 Calabi--Yau 流形的环面退化出发，第一步是构造对偶相交复形。然后取它的（离散）Legendre 变换，并试图从这个 Legendre 对偶重建 Calabi--Yau 流形的镜像环面退化。这样一来，他们就可以完全忘掉 SYZ 纤维化。其主张是：所有相关信息都编码在对偶相交复形的热带几何中，而后者是一个带奇点的整仿射流形，扮演 SYZ 纤维化之底的角色。
% [Tutor Note: "dual intersection complex"（对偶相交复形）是 Gross--Siebert 纲领的核心对象：它是退化中心纤维的相交复形的对偶，带有整仿射结构，充当 SYZ 纤维化之底的角色。]

利用上述 Kontsevich--Soibelman 的关键引理，再加上对数结构的使用以及热带几何技巧等许多新想法，Gross 与 Siebert 最终成功地给出了任意维数下重建问题的解答 [87]。给定一个满足某些假设、并带有诸如多面体分解等附加结构的带奇点整仿射流形，他们构造出一个 Calabi--Yau 流形的环面退化，它可以通过 $B$ 中的热带树显式且典范地描述。他们构造的一个重要特点是：所构造的 Calabi--Yau 流形定义在 $\C$ 上，而不是刚性解析空间。
% [Tutor Note: "defined over $\C$"（定义在 $\C$ 上）与函数域上的刚性解析空间相对；Gross--Siebert 的构造比 Kontsevich--Soibelman 的非阿基米德解析构造更"代数"。]

现在的目标是通过热带世界来获得对镜像对称更概念化的理解。在 B-侧，Gross 与 Siebert 猜想其构造中的形变参数是一个典范坐标\footnote{这一点最近由 Ruddat--Siebert [155] 所确认。}，并且 Calabi--Yau 流形族的周期积分可以用 $B$ 中的热带圆盘来表达。最先的证据来自局部情形，例如 [87, 注 5.1] 中的局部 $\PP^2$ 例子。

在 A-侧，人们希望用环面退化中心纤维的 Gromov--Witten 理论来理解光滑纤维的 Gromov--Witten 理论。最近 Abramovich--Chen [34, 4] 与 Gross--Siebert [89] 的独立工作发展了对数 Gromov--Witten 不变量的理论，推广了 Li--Ruan [123]、Ionel--Parker [101, 102] 与 Jun Li [124] 先前关于相对 Gromov--Witten 理论的工作。这是朝着最终目标迈出的重要一步。如果还能进一步证明热带曲线/圆盘与全纯曲线/圆盘之间的一般对应定理（沿 Mikhalkin [138, 139]、Nishinou--Siebert [147] 与 Nishinou [146, 145] 的工作的思路），那么我们就能把一个 Calabi--Yau 流形的 A-侧（即 Gromov--Witten 理论）与热带世界联系起来。

尽管还有许多工作要做，这已经勾勒出一幅有前景的美丽图景，用热带几何来解释镜像对称的几何。我们建议读者参阅该纲领的创立者所写的综述文章 [88, 75]，以了解 Gross--Siebert 纲领的概貌。

% [原文 p.9]
\section{非紧 Calabi--Yau 流形的 SYZ}
缺乏特殊拉格朗日环面纤维化的例子，是把原始 SYZ 提议实施于紧 Calabi--Yau 流形的主要障碍之一（这或许也是 Gross 与 Siebert 试图发展代数几何版本的主要原因之一）。但在非紧情形有大量例子可以找到显式的特殊拉格朗日环面纤维化，例如 Goldstein [65] 与 Gross [71] 在环面 Calabi--Yau 流形情形所构造的那些。

此外，计数从开 Riemann 面到流形的稳定映射的开 Gromov--Witten 不变量，在环面情形由 Fukaya--Oh--Ohta--Ono 的工作 [57, 58, 59] 良定义（更一般地，在 $S^1$-等变情形由 Liu 的工作 [130] 良定义）。因此，直接在环面 Calabi--Yau 流形上实施 SYZ 提议是完全合理的，无需退回到热带世界与代数几何版本。

这把我们带入局部镜像对称的领域；局部镜像对称由紧 Calabi--Yau 流形的镜像对称在 Kähler 模空间与复模空间中取某些极限而得到 [107]。由于这种镜像对称提供了许多有趣的例子并有大量应用，它吸引了物理学家与数学家的广泛关注 [121, 38, 97, 71, 72, 109, 66, 99, 100, 50, 51, 110, 156]。
% [Tutor Note: "local mirror symmetry"（局部镜像对称）不是局部化的镜像对称，而是由紧 Calabi--Yau 情形在 Kähler 与复模空间取极限得到的、适用于非紧环面 Calabi--Yau 流形的镜像对称。]

设 $X$ 是一个 $n$ 维环面 Calabi--Yau 流形（它必然是非紧的）。为了实施 SYZ 构造，我们使用 Goldstein 与 Gross 构造的特殊拉格朗日环面纤维化 $\mu:X\to B$；这样的纤维化是非环面的，意思是它不是 $X$ 上 Hamiltonian $T^n$-作用所对应的通常矩映射。这个 SYZ 纤维化的判别轨迹已被 Gross 详细分析，可以被显式描述。

在拓扑上，底 $B$ 就是 $\R^n$ 中的一个上半空间，它是一个同时带奇点与边界的整仿射流形。边界 $\partial B\subset B$ 的原像是 $X$ 中一个非环面光滑超曲面 $D$（即环面主除子并集的一个光滑化）。判别轨迹 $\Gamma\subset B$ 是一个实余维二的热带子簇，包含在超平面 $H$ 中，我们称 $H$ 为 $B$ 中的墙。按定义，SYZ 纤维化底内的墙是那些在 $X$ 中界定 Maslov 指标 0 全纯圆盘的拉格朗日环面纤维所在的轨迹。它把底分成不同的腔室，拉格朗日环面纤维在这些腔室上的 Floer 理论行为不同。在 Gross 纤维化的情形，墙 $H\subset B$ 平行于边界超平面 $\partial B$，把底分成两个腔室 $B_+$ 与 $B_-$。
% [Tutor Note: 墙上的纤维会界定 Maslov 指标 0 的"常值"全纯圆盘，因而其 Floer 理论（及开 Gromov--Witten 不变量）在穿过墙时发生跳变；这解释了为何必须按腔室分别取纤维对偶再做墙穿越粘接。]

我们考虑 $X$ 中由 $\mu$ 在 $B_+$ 与 $B_-$ 上的纤维所界定的 Maslov 指标 2 全纯圆盘的（虚拟）计数；这些圆盘与超曲面 $D$ 相交于一点，重数为 1。当一点从 $B_-$ 穿过墙移动到 $B_+$ 时，由相应拉格朗日环面纤维界定的 Maslov 指标 2 圆盘的虚拟数目------更确切地说，亏格 0 开 Gromov--Witten 不变量------发生跳变，表现出墙穿越现象。

这一点已由 Auroux [8, 9] 与 Chan--Lau--Leung [20] 利用 Fukaya--Oh--Ohta--Ono [56] 发展的精密工具进行了分析。注意在这种情形没有散射现象，因为只有一个墙。通过在底的每个腔室上取纤维对偶，再按照墙穿越公式把所得的各块粘接起来，我们得到瞬子修正的（或 SYZ 的）镜像 $\check X$，它是 $X$ 的（复化）Kähler 模空间上的一族仿射超曲面（它们也是非紧 Calabi--Yau 流形）[20, 3]；这与物理学家的预言一致 [38, 97]。
% [Tutor Note: 散射图只在多于一个墙时出现；Gross 纤维化只有一个墙 $H$，故此处只有墙穿越、没有散射，与 p.8 的散射图讨论形成对照。]

% [原文 p.10]
这个 SYZ 构造非常精确，因为它确切地告诉我们：对 $X$ 上任给一个辛结构，镜像 $\check X$ 上与之对应的复结构是哪一个------镜像 $\check X$ 的定义方程是一个显式表达式，完全用 $X$ 的 Kähler 参数与圆盘计数不变量写出。例如，$X=K_{\PP^2}$ 的 SYZ 镜像由下式给出\footnote{更确切地说，$K_{\PP^2}$ 的 SYZ 镜像由 Landau--Ginzburg 模型 $(\check X,W=u)$ 给出；见下一节。}：
\begin{equation}
\check X=\Bigl\{(u,v,z_1,z_2)\in\C^2\times(\C^\times)^2\ \Big|\ uv=\delta(q)+z_1+z_2+\frac{q}{z_1z_2}\Bigr\},
\tag{2.6.1}
\end{equation}
其中 $q$ 是 Kähler 参数，它度量 $K_{\PP^2}$ 在 $\PP^2$ 上零截面内一条射影直线的辛面积，而
\begin{equation}
\delta(q)=\sum_{k=0}^{\infty}n_kq^k
\tag{2.6.2}
\end{equation}
是亏格 0 开 Gromov--Witten 不变量的生成级数。
% [Tutor Note: $K_{\PP^2}$ 表示 $\PP^2$ 上典范丛 $K_{\PP^2}=\mathcal{O}_{\PP^2}(-3)$ 的全空间，它是非紧 Calabi--Yau 曲面（复二维）；不要读成"域 $K$ 上的 $\PP^2$"。]
% [Tutor Note: (6.1) 中镜像 $\check X$ 是 $\C^2\times(\C^\times)^2$ 中由 $uv=\delta(q)+z_1+z_2+q/(z_1z_2)$ 定义的仿射超曲面；$u,v$ 是 $\C^2$ 的坐标，$z_1,z_2$ 是 $(\C^\times)^2$ 的坐标，两者不可混同。]

此外，SYZ 构造自然地定义了 SYZ 映射，它是从 $X$ 的 Kähler 模空间到 $\check X$ 的复模空间的映射。正如 Gross 与 Siebert 所猜想的那样 [87, 猜想 0.2 与注 5.1]，SYZ 镜像族应当用典范坐标写出。在环面 Calabi--Yau 情形，这等价于说 SYZ 映射是镜像映射的逆。

关于这个猜想在环面 Calabi--Yau 曲面与三维流形情形的证据在 [20, 117] 中给出；Chan--Lau--Tseng [21] 在 $X$ 是紧环面 Fano 流形上典范丛全空间的情形证明了该猜想。最近，通过应用 orbifold 技巧，[19] 对所有环面 Calabi--Yau 流形（以及 orbifold）证明了该猜想。我们称之为开镜像定理，因为它为（逆）镜像映射提供了枚举意义，并给出计算所有亏格 0 开 Gromov--Witten 不变量的有效方法。\footnote{在他们的 ICM 报告 [82] 中，Gross 与 Siebert 勾勒了该猜想的另一个证明，使用对数 Gromov--Witten 理论 [34, 4, 89]，并把全纯圆盘换成了热带圆盘。最近，Lau [116] 应用 [19] 中的结果证明了开 Gromov--Witten 不变量的生成函数（例如 $X=K_{\PP^2}$ 情形中的 $\delta(q)$）是 Gross--Siebert 意义下的 slab 函数。}

证明这些结果的主要挑战在于计算 Fukaya--Oh--Ohta--Ono [57] 定义的亏格 0 开 Gromov--Witten 不变量。由于全纯圆盘的模空间通常是高度障碍的，这些不变量一般很难计算。目前只有很少的技巧可用（例如开/闭相等式、环面镜像定理、退化技巧等）。例如，(2.6.2) 中的不变量可以算得为
\[
n_k=1,-2,5,-32,286,-3038,35870,\dots
\]
其中 $k=0,1,2,3,4,5,6,\dots$，这与 [66] 中的周期计算一致。
% [Tutor Note: 这里 $n_k$ 是亏格 0 开 Gromov--Witten 不变量（圆盘计数），不是陈数或 Betti 数；其数值 $1,-2,5,-32,286,-3038,35870,\dots$ 是原文给出的、可用周期计算核对的序列。]

我们应当指出，SYZ 构造也可以在相反方向上进行 [3]（亦见 [18, 第 5 节]）。
例如，从式 (2.6.1) 给出的二次曲线丛（conic bundle）出发，可以用与 [65, 71] 类似的技术构造一个 SYZ 纤维化。
